\documentclass[11pt, a4paper]{article}

\usepackage[utf8]{inputenc}
\usepackage[T1]{fontenc}
\usepackage[french]{babel}
\usepackage{amsmath, amssymb}
\usepackage{geometry}
\usepackage{graphicx}
\usepackage{tabularx}
\usepackage{booktabs}
\usepackage{xcolor}
\usepackage{listings}
\usepackage{tikz}
\usepackage{hyperref}

% Configuration des marges
\geometry{
    top=2cm,
    bottom=2cm,
    left=1.5cm,
    right=1.5cm
}

% Configuration pour les blocs de code VHDL
\lstset{
    basicstyle=\ttfamily\small,
    breaklines=true,
    frame=single,
    tabsize=4,
    showstringspaces=false,
    keywordstyle=\color{blue}\bfseries,
    commentstyle=\color{green!50!black},
    stringstyle=\color{red},
    numbers=left,
    numberstyle=\tiny\color{gray}
}

% Commande pour les placeholders d'images
\newcommand{\imageplaceholder}[2]{
    \begin{center}
        \fbox{\begin{minipage}{#1}\centering\vspace{1cm}\textit{#2}\vspace{1cm}\end{minipage}}
    \end{center}
}

% Commande pour afficher la correction avec explication
\newcommand{\corr}[1]{
    \vspace{0.3cm}
    \begin{center}
    \fcolorbox{blue}{blue!5}{
    \begin{minipage}{0.95\textwidth}
    \color{blue} \textbf{Correction \& Raisonnement :} \\
    #1
    \end{minipage}}
    \end{center}
    \vspace{0.3cm}
}

\begin{document}

% ==================== EN-TÊTE ====================
\noindent
\begin{tabularx}{\textwidth}{@{}X c r@{}}
\textbf{SN361} & \textbf{Session 1 - Année 2017-2018} & \textbf{Documents selon sujet} \\
\textbf{VHDL} & \textbf{Examen} & \textbf{Calculatrice selon sujet} \\
\end{tabularx}
\vspace{0.2cm}
\hrule
\vspace{0.4cm}

% ==================== PROBLÈME I ====================
\begin{center}
    \textbf{\Large Problème I - Conception d'un registre - 40 pts}
\end{center}
\vspace{-0.2cm}
\hrule
\vspace{0.3cm}

\noindent \textbf{Question 1.1} : Dessinez en utilisant uniquement des bascules DFF (possédant des entrées reset et enable) le circuit d'un registre $n$ bits à chargement parallèle et sortie parallèle capturant de l'entrée \texttt{Din} sur la sortie \texttt{Dout} sur les fronts montants de l'horloge \texttt{clk} avec :
\begin{itemize}
    \item une entrée \texttt{enable} active à l'état haut autorisant le chargement de l'entrée \texttt{Din}.
    \item une entrée de données \texttt{Din} sur $n$ bits.
    \item une sortie \texttt{Dout} sur $n$ bits.
    \item une remise à zéro \texttt{Reset} synchrone active à l'état haut.
\end{itemize}

\corr{
\textbf{Étape 1 : Structure globale du registre.}\\
Un registre sur $n$ bits n'est rien d'autre qu'une mise en parallèle de $n$ bascules de type D indépendantes, permettant de mémoriser simultanément un mot complet de $n$ bits.

\textbf{Étape 2 : Câblage des signaux de contrôle.}\\
Tous les signaux de commande ou de rythme (\texttt{clk}, \texttt{reset}, \texttt{enable}) sont connectés de manière commune et distribués à toutes les bascules du circuit en parallèle.

\textbf{Étape 3 : Câblage des bus de données.}\\
Le bit de poids $i$ du bus d'entrée (\texttt{Din[i]}) attaque directement l'entrée \texttt{D} de la bascule de rang $i$. Sa sortie \texttt{Q} correspond directement au bit $i$ du bus de sortie (\texttt{Dout[i]}). Ainsi, sur un front montant, si \texttt{enable=1} et \texttt{reset=0}, chaque bascule recopie l'état de son bit d'entrée respectif.
}

\noindent \textbf{Question 1.2} : Ecrire en VHDL RTL l'entité générique (\texttt{registerN}) et l'architecture (\texttt{rtl}) du composant précédent [10].

\corr{
\textbf{Étape 1 : Définition de l'entité.}\\
On déclare l'entité avec un paramètre générique pour la taille $n$ et les ports d'entrée/sortie sur forme de vecteurs.

\textbf{Étape 2 : Implémentation du comportement RTL.}\\
On utilise un processus séquentiel sensible à l'horloge. Le \texttt{reset} étant synchrone, il est testé à l'intérieur de la condition \texttt{rising\_edge(clk)}, avec une priorité supérieure à \texttt{enable}.
\begin{lstlisting}[language=VHDL]
library IEEE;
use IEEE.std_logic_1164.all;

entity registerN is
    generic (
        n : integer := 8
    );
    port (
        clk    : in  std_logic;
        reset  : in  std_logic;
        enable : in  std_logic;
        Din    : in  std_logic_vector(n-1 downto 0);
        Dout   : out std_logic_vector(n-1 downto 0)
    );
end entity registerN;

architecture rtl of registerN is
begin
    process(clk)
    begin
        if rising_edge(clk) then
            if reset = '1' then
                Dout <= (others => '0'); -- Remise a zero synchrone
            elsif enable = '1' then
                Dout <= Din;             -- Chargement parallele
            end if;
        end if;
    end process;
end architecture rtl;
\end{lstlisting}
}

\noindent \textbf{Question 1.3} : Combien de bascules seront générées par la synthèse de cette description rtl [5].

\corr{
\textbf{Étape 1 : Règles d'inférence de bascules.}\\
La synthèse logique infère systématiquement une bascule D pour chaque bit d'un signal assigné de façon conditionnelle sous l'évaluation d'un front d'horloge.

\textbf{Étape 2 : Analyse et dénombrement.}\\
Le processus mémorise la sortie \texttt{Dout}. Comme cette sortie est un vecteur défini sur la plage \texttt{(n-1 downto 0)}, elle comporte physiquement $n$ bits indépendants.
$$ \mathbf{Total = n \text{ bascules}} $$
}

\noindent \textbf{Question 1.4} : Ecrire en VHDL structurel l'architecture (\texttt{struct}) de la même entité (\texttt{registerN}) précédente [15]. On dispose du composant suivant :
\begin{lstlisting}[language=VHDL]
entity bascule is
port(d, clk, reset, enable: in std_logic;
     q: out std_logic);
end entity bascule;
\end{lstlisting}

\corr{
\textbf{Étape 1 : Déclaration du composant.}\\
Il faut d'abord déclarer le composant unitaire \texttt{bascule} dans l'en-tête de l'architecture pour qu'il soit connu du compilateur.

\textbf{Étape 2 : Instanciation répétée via une boucle.}\\
Puisque le nombre de bascules dépend d'un paramètre générique ($n$), on utilise la syntaxe \texttt{for ... generate} pour instancier automatiquement les $n$ composants matériels en mappant le bit $i$ de chaque vecteur à l'instance de rang $i$.
\begin{lstlisting}[language=VHDL]
architecture struct of registerN is
    component bascule
        port(d, clk, reset, enable: in std_logic;
             q: out std_logic);
    end component;
begin
    gen_bascules: for i in 0 to n-1 generate
        bascule_inst : bascule
            port map (
                d      => Din(i),
                clk    => clk,
                reset  => reset,
                enable => enable,
                q      => Dout(i)
            );
    end generate gen_bascules;
end architecture struct;
\end{lstlisting}
}

\newpage

% ==================== PROBLÈME II ====================
\begin{center}
    \textbf{\Large Problème II - Conception d'un « compteur générique » - 50 pts}
\end{center}
\vspace{-0.2cm}
\hrule
\vspace{0.3cm}

\noindent \textbf{Question 2.1} : Ecrire en VHDL RTL l'entité (\texttt{compteur\_moduloN}) avec un paramètre générique \texttt{N} et l'architecture (\texttt{rtl}) d'un compteur comptant sur 4 bits et modulo N [15].
\begin{itemize}
    \item Entrées : horloge (\texttt{clk}) sur front montant, reset (\texttt{rst}) asynchrone actif haut, enable (\texttt{en}) actif haut.
    \item Sorties : valeur (\texttt{count}, sur 4 bits), signal (\texttt{fin}) qui passe à '1' quand la valeur max (N-1) est atteinte.
\end{itemize}

\corr{
\textbf{Étape 1 : Définition du comportement (Modulo).}\\
Un compteur modulo $N$ s'incrémente de 0 jusqu'à la borne $N-1$, puis retourne à 0 au cycle suivant.

\textbf{Étape 2 : Processus de comptage (Reset asynchrone).}\\
Puisque le signal \texttt{rst} est défini comme asynchrone, il doit être inclus dans la liste de sensibilité du processus et doit être évalué en priorité de manière isolée, sans attendre le front d'horloge.

\textbf{Étape 3 : Logique de sortie.}\\
La valeur courante \texttt{count} est recopiée depuis un signal interne de type \texttt{unsigned}. Le drapeau \texttt{fin} est un signal combinatoire qui compare directement la valeur du compteur avec la borne.
\begin{lstlisting}[language=VHDL]
library IEEE;
use IEEE.std_logic_1164.all;
use IEEE.numeric_std.all;

entity compteur_moduloN is
    generic ( N : integer := 10 );
    port (
        clk   : in  std_logic;
        rst   : in  std_logic;
        en    : in  std_logic;
        count : out std_logic_vector(3 downto 0);
        fin   : out std_logic
    );
end entity compteur_moduloN;

architecture rtl of compteur_moduloN is
    signal count_int : unsigned(3 downto 0);
begin
    process(clk, rst)
    begin
        if rst = '1' then
            count_int <= (others => '0');   -- Reset asynchrone
        elsif rising_edge(clk) then
            if en = '1' then
                if count_int = N - 1 then
                    count_int <= (others => '0'); -- Bouclage modulo N
                else
                    count_int <= count_int + 1;   -- Incrementation
                end if;
            end if;
        end if;
    end process;

    count <= std_logic_vector(count_int);
    fin <= '1' when (count_int = N - 1) else '0';
end architecture rtl;
\end{lstlisting}
}

\noindent \textbf{Question 2.2} : Combien de bascules seront générées par la synthèse logique de votre description précédente [5] ?

\corr{
\textbf{Étape 1 : Analyse des signaux mémorisés sous horloge.}\\
Le seul signal assigné dans la section synchrone du processus (sous la condition du front d'horloge) est le signal interne \texttt{count\_int}.

\textbf{Étape 2 : Calcul de la taille.}\\
L'énoncé contraint le compteur à produire une sortie sur 4 bits (\texttt{3 downto 0}), indépendamment de la valeur de N. Le signal \texttt{fin} quant à lui est purement combinatoire.
$$ \mathbf{Total = 4 \text{ bascules}} $$
}

\noindent \textbf{Question 2.3} : Ecrire en VHDL un testbench permettant de simuler l'architecture précédente pour un modulo $N=8$ [15].

\corr{
\textbf{Étape 1 : Structure globale du Testbench.}\\
L'entité est sans port. Dans l'architecture, on déclare les signaux de stimulation et on instancie l'unité sous test (UUT) en fixant \texttt{N = 8} via un `generic map`.

\textbf{Étape 2 : Génération de la porteuse.}\\
On crée un processus autonome qui fait basculer le signal \texttt{clk} en permanence pour créer la fréquence.

\textbf{Étape 3 : Scénario séquentiel.}\\
Un processus de stimuli applique le \texttt{rst}, libère le composant, et active le \texttt{en} assez longtemps pour observer au moins un tour complet.
\begin{lstlisting}[language=VHDL]
library IEEE;
use IEEE.std_logic_1164.all;

entity tb_compteur is
end entity tb_compteur;

architecture behavior of tb_compteur is
    signal clk, rst, en, fin : std_logic := '0';
    signal count : std_logic_vector(3 downto 0);
    constant clk_period : time := 10 ns;
begin
    uut: entity work.compteur_moduloN
        generic map ( N => 8 )
        port map ( clk => clk, rst => rst, en => en, count => count, fin => fin );

    clk_process : process begin
        clk <= '0'; wait for clk_period/2;
        clk <= '1'; wait for clk_period/2;
    end process;

    stim_proc: process begin
        rst <= '1'; en <= '0'; wait for 20 ns;
        
        rst <= '0'; en <= '1'; -- Demarrage
        wait for 100 ns;       -- Attente de 10 cycles (modulo 8 sera visible)
        
        en <= '0';             -- Arret
        wait;
    end process;
end architecture behavior;
\end{lstlisting}
}

\noindent \textbf{Question 2.4} : Dessinez les chronogrammes des signaux de l'entité correspondant très exactement à la simulation que donnerait votre testbench précédent [15]. Indiquez clairement l'échelle des temps sur l'axe temporel.

\corr{
\textbf{Étape 1 : Période de maintien sous reset (0 à 20 ns).}\\
Le signal \texttt{rst} force le maintien du compteur \texttt{count} à "0000" et \texttt{fin} à '0'. L'horloge bat toutes les $10\text{ ns}$ mais sans effet sur la sortie.

\textbf{Étape 2 : Phase d'incrémentation (20 ns à 100 ns).}\\
Dès le relâchement du reset et l'activation de \texttt{en}, la valeur s'incrémente au rythme des fronts montants :
$0 \rightarrow 1 \rightarrow 2 \rightarrow 3 \rightarrow 4 \rightarrow 5 \rightarrow 6 \rightarrow 7$.
Lors du passage à la valeur 7 (à $t = 85\text{ ns}$), le signal \texttt{fin} passe immédiatement à '1' de façon combinatoire car il vaut $N-1$.

\textbf{Étape 3 : Bouclage modulo (100 ns).}\\
Au front d'horloge suivant ($t = 95\text{ ns}$), le compteur ne peut pas passer à 8. Il reboucle à "0000" (condition modulo vérifiée). Le signal \texttt{fin} retombe instantanément à '0'. Le testbench fige l'exécution ensuite.
}

\vfill
% ==================== PIED DE PAGE ====================
\hrule
\vspace{0.1cm}
\noindent
Équipe Pédagogique \hfill Esisar \hfill Page 1/1

\end{document}